\chapter{Experimental Setup and Results}
\label{chap:experiment}

\section{Benchmark Environments}
The empirical evaluation is conducted on five sparse-reward MiniGrid environments \cite{chevalier2018minigrid}. These environments are chosen because they represent different forms of exploration difficulty while maintaining a common gridworld interface.

\begin{table}[H]
\centering
\caption{MiniGrid environments used in the experiments}
\begin{tabular}{|m{4cm}|m{9cm}|}
\hline
\textbf{Environment} & \textbf{Purpose in the evaluation} \\ \hline
DoorKey-8x8 & Tests whether the agent can discover a key, unlock a door, and shift from exploration to goal-directed behavior. \\ \hline
Empty-16x16 & Represents an extreme sparse-reward setting with almost no intermediate feedback before the goal is reached. \\ \hline
RedBlueDoors-8x8 & Requires opening two doors in the correct order, testing whether the method can suppress irrelevant exploration. \\ \hline
UnlockPickup & Extends DoorKey with an additional object manipulation stage, creating a hierarchical sparse-reward task. \\ \hline
KeyCorridorS3R3 & A longer-horizon task with multiple rooms, keys, and locked doors, used to assess sustained exploration ability. \\ \hline
\end{tabular}
\end{table}

These tasks share several useful properties: sparse terminal rewards, egocentric partial observations, step penalties that discourage random motion, and stochastic initialization over multiple random seeds. Together, they form a suitable test bed for evaluating whether adaptive intrinsic reward scaling improves both efficiency and robustness.

\section{Baselines and Evaluation Criteria}
The proposed ACWI method is compared against two categories of baselines:
\begin{itemize}
    \item PPO trained only with extrinsic rewards;
    \item PPO combined with ICM using fixed intrinsic coefficients $\beta \in \{0.1, 0.2, 0.5, 1, 2\}$.
\end{itemize}

This comparison is important because it isolates the contribution of adaptive scaling from the contribution of intrinsic motivation itself. PPO provides a lower bound on exploration ability, while fixed-coefficient PPO+ICM shows how sensitive performance is to manual tuning.

The evaluation focuses on four aspects:
\begin{itemize}
    \item episode return over training time;
    \item sample efficiency in the early and middle stages of learning;
    \item variance across random seeds;
    \item qualitative behavior of the learned scaling coefficient $\beta(s)$.
\end{itemize}

\section{Experimental Configuration}
All methods share the same PPO backbone and ICM architecture so that the comparison is fair. The adapted paper fixes the global intrinsic coefficient at $\alpha = 0.001$ and clips the learned Beta Network output to $[0.1, 2.0]$. The Beta Network uses a hidden representation size of 256 with encoder depth 2. The regularization weight for the correlation objective is $\lambda_{\text{reg}} = 10^{-3}$, the reference value is $\beta_0 = 1.0$, gradient clipping is set to 1.0, and weight decay is $10^{-6}$.

\begin{table}[H]
\centering
\caption{Key hyperparameters reported for ACWI}
\begin{tabular}{|m{6cm}|m{4cm}|}
\hline
\textbf{Hyperparameter} & \textbf{Value} \\ \hline
Intrinsic strength $\alpha$ & 0.001 \\ \hline
Beta output bounds & $[0.1, 2.0]$ \\ \hline
Encoder size & 256 \\ \hline
Encoder depth & 2 \\ \hline
Regularization weight $\lambda_{\text{reg}}$ & $10^{-3}$ \\ \hline
Reference value $\beta_0$ & 1.0 \\ \hline
Gradient clipping & 1.0 \\ \hline
Weight decay & $10^{-6}$ \\ \hline
\end{tabular}
\end{table}

The results are reported over five random seeds. This setup is sufficient to assess whether adaptive scaling consistently improves learning behavior rather than benefiting from a single favorable initialization.

\section{Results}
The experiments show several consistent trends. First, ACWI performs strongly across environments with sparse but still informative extrinsic rewards. In DoorKey-8x8, RedBlueDoors-8x8, UnlockPickup, and KeyCorridorS3R3, the adaptive approach achieves faster learning and lower variance than most fixed-coefficient baselines. This result indicates that ACWI improves sample efficiency and makes exploration behavior more stable across runs.

Second, the comparison against fixed-coefficient PPO+ICM highlights the difficulty of manual tuning. A coefficient that works well in one environment can be too small or too large in another. Small values underemphasize exploration and slow down learning, while large values can inject excessive noise and interfere with exploitation. ACWI mitigates both problems by learning to amplify intrinsic rewards only in states where they correlate with future extrinsic success.

Third, the method exhibits a natural transition from exploration to exploitation. As training progresses and the policy begins to discover task-relevant behaviors, the effective intrinsic contribution tends to decrease. This suggests that ACWI does not merely add curiosity-driven noise, but instead learns to reduce exploration pressure once the task structure becomes clearer.

\section{Analysis of the Learned Scaling Behavior}
The paper reports that the learned distribution of $\beta(s)$ evolves differently across environments. In DoorKey-8x8 and RedBlueDoors-8x8, the distribution is initially concentrated near its reference value, then gradually broadens and becomes more structured. This pattern suggests that the Beta Network learns to partition the state space into regions where different exploration intensities are appropriate. Toward the end of training, the distribution shifts toward lower values, reflecting the decreasing marginal value of intrinsic rewards once the policy has matured.

In contrast, Empty-16x16 behaves differently. Because extrinsic rewards remain almost zero for most trajectories, the correlation objective receives very weak supervision. As a result, the Beta Network stays close to its initialization and exhibits limited adaptation. Importantly, this failure mode is graceful rather than unstable: the method effectively defaults to a nearly fixed intrinsic coefficient instead of diverging. This observation confirms both a strength and a limitation of ACWI. The method is robust when correlation signals are weak, but meaningful adaptation requires at least occasional extrinsic feedback.

\section{Discussion}
Overall, the empirical evidence supports the main claim of the thesis: state-dependent intrinsic reward scaling can improve sparse-reward reinforcement learning when the environment provides enough structure for the agent to relate exploration to eventual success. ACWI is especially beneficial in environments where exploration needs change across states and over time. The method is less effective in extremely sparse tasks where the extrinsic return is too weak to supervise adaptation. Therefore, ACWI should be viewed as a practical middle ground between fixed intrinsic scaling and more expensive meta-learning approaches.
